\subsection{指数映射}

定理 4. 设 \(\mathbf{G}\) 是李群。存在唯一一个 \(\mathbf{G}\) 的指数映射，定义在 \(\mathrm{L}(\mathbf{G})\) 上。

存在一个凸开邻域 \(\mathbf{U}\)，它位于 \(\mathrm{L}(\mathbf{G})\) 中并包含 0，以及一个 \(\mathbf{G}\) 的指数映射，定义在 \(\mathbf{U}\) 上。通过把 \(\mathbf{U}\) 取得充分小，可以假设

\[
\phi ((\lambda + \lambda^ {\prime}) a) = \phi (\lambda a) \phi (\lambda^ {\prime} a)\tag{1}
\]

其中 \(a \in \mathrm{L}(\mathbf{G})\)，\(\lambda, \lambda'\) 属于 \(\mathbf{K}\)，且 \(\lambda a, \lambda'a\)、\((\lambda + \lambda')a\) 属于 \(\mathbf{U}\)。

令 \(a \in \mathrm{L}(\mathrm{G})\)。存在一个整数 \(n > 0\)，使 \(\frac{1}{n} a \in \mathrm{U}\)。若 \(m\) 是满足 \(\frac{1}{m} a \in \mathrm{U}\) 的另一个整数，则 \(\frac{1}{nm} a \in \mathrm{U}\)，且关系 (1) 蕴含

\[
\phi \left(\frac {1}{n} a\right) = \left(\phi \left(\frac {1}{n m} a\right)\right) ^ {m}, \quad \phi \left(\frac {1}{m} a\right) = \left(\phi \left(\frac {1}{n m} a\right)\right) ^ {n};
\]

从而 \(\left(\phi\left(\frac{1}{n}a\right)\right)^n = \left(\phi\left(\frac{1}{m}a\right)\right)^m\)。存在一个延拓 \(\psi: \mathrm{L}(G) \to G\)，它是 \(\phi\) 的延拓，并满足 \(\psi(a) = \left(\phi\left(\frac{1}{n}a\right)\right)^n\) 对 \(a \in \mathrm{L}(G)\) 以及 \(n\) 满足 \(>0\) 且 \(\frac{1}{n}a \in \mathrm{U}\) 时成立。显然 \(\psi\) 是解析的，也是 \(G\) 的指数映射。若 \(\psi': \mathrm{L}(G) \to G\) 是 \(G\) 的一个指数映射，则 \(\psi\) 和 \(\psi'\) 在 0 的某个邻域上重合；由于 \(\mathrm{L}(G)\) 连通，故二者相等。

今后提到 G 的指数映射时，均指定理 4 中所考虑的映射。若不会引起混淆，将其记为 \(\exp_{G}\) 或 exp。

例。设 A 是完备赋范含幺结合代数。则 \(\exp_{\mathbf{A}}\) 是 Chapter II, § 7, no. 3 中定义的指数映射。

命题 7. 设 G 是李群，a 是 L(G) 的元素。从 K 到 G 的映射 \(\lambda \mapsto \exp(\lambda a)\) 是唯一的李群 K 到 G 的态射 \(\phi\)，满足 \((\mathrm{T}_0\phi)1 = a\)。

映射 \((\lambda, \lambda') \mapsto \exp(\lambda a) \exp(\lambda' a)\) 和 \((\lambda, \lambda') \mapsto \exp(\lambda + \lambda') a\) 从 \(K \times K\) 到 \(G\)，都是解析的，并在 \((0, 0)\) 的某个邻域上重合。由于 \(K \times K\) 连通，这两个映射相等。因此 \(\phi: \lambda \mapsto \exp(\lambda a)\) 是从 \(K\) 到 \(G\) 的李群态射。从 \(0\) 处到映射 \(\lambda \mapsto \lambda a\) 的切映射就是映射 \(\lambda \mapsto \lambda a\)；且 \(T_0(\exp) = \mathrm{Id}_{\mathrm{L}(\mathrm{G})}\)；所以 \((T_0\phi)1 = a\)。命题的唯一性断言由定理 1 得出。

命题 8. 设 G 是李群。对于 L(G) 中所有 x, y 以及整数 n，

\begin{enumerate}
\def\labelenumi{(\arabic{enumi})}
\setcounter{enumi}{1}
\tightlist
\item
\end{enumerate}

\[
\exp (x + y) = \lim _ {n \rightarrow + \infty} \left(\left(\exp \frac {1}{n} x\right)\left(\exp \frac {1}{n} y\right)\right) ^ {n}\tag{2}
\]

\[
\exp [ x, y ] = \lim _ {n \rightarrow + \infty} \left(\left(\exp \frac {1}{n} x\right)\left(\exp \frac {1}{n} y\right)\left(\exp \frac {1}{n} x\right) ^ {- 1} \left(\exp \frac {1}{n} y\right) ^ {- 1}\right) ^ {n ^ {2}}\tag{3}.
\]

由命题 7，取 \(\lambda = \frac{1}{n}\)，这由 § 4, no. 3 的 Proposition 4 得出。

命题 9. 设 G 是复李群，G' 是 G 的底层实李群。则 \(\exp_{G} = \exp_{G'}\)。

这由 §4, no. 3 的 Proposition 5 以及 \(\exp_{G}\) 和 \(\exp_{G'}\) 的解析性得出。

命题 10. 设 G 和 H 是李群，\(\phi\) 是从 G 到 H 的态射。

\begin{enumerate}
\def\labelenumi{(\roman{enumi})}
\item
  \(\phi \circ \exp_{\mathrm{G}} = \exp_{\mathrm{H}} \circ \mathrm{L}(\phi)\) .
\item
  若 G 是 H 的积分子群，则 \(\exp_{G} = \exp_{H} | L(G)\)。
\end{enumerate}

等式 (i) 两边都是从 L(G) 到 H 的解析映射，并在 0 的某个邻域上重合（§ 4, Proposition 8, no. 4），因而相等。断言 (ii) 是 (i) 的特殊情形。

推论 1. 设 G 是李群，G' 是 G 的李子群，且 \(a \in \mathrm{L}(\mathrm{G})\)。以下条件等价：

\begin{enumerate}
\def\labelenumi{(\roman{enumi})}
\item
  \(a \in \mathbf{L}(\mathbf{G}')\) ;
\item
  当 \(\exp (\lambda a)\in \mathbf{G}'\) 对所有 \(\lambda \in \mathbf{K}\) 且 \(|\lambda |\) 充分小时成立；
\item
  \(\exp(\lambda a) \in \mathbf{G}'\) 对所有 \(\lambda \in K\) 成立。
\end{enumerate}

论证与 § 4, no. 4, Corollary 1 to Proposition 8 中的相同。

推论 2. 设 \(\mathbf{G}\) 是李群，\(\mathbf{H}\) 是 \(\mathbf{G}\) 的积分子群，且 \(a \in \mathbf{L}(\mathbf{G})\)。考虑以下条件：

\begin{enumerate}
\def\labelenumi{(\roman{enumi})}
\item
  \(a \in \mathrm{L}(\mathrm{H})\) ;
\item
  \(\exp_{\mathbf{G}}(\lambda a)\in \mathbf{H}\) 对所有 \(\lambda \in \mathbf{K}\) 成立。
\end{enumerate}

则 (i) \(\Rightarrow\) (ii)。若 H 的拓扑有可数基，则 (i) \(\Leftrightarrow\) (ii)。

令 \(i\) 为从 \(\mathbf{H}\) 到 \(\mathbf{G}\) 的规范嵌入。若 \(a \in \mathbf{L}(\mathbf{H})\)，则

\[
\exp_ {\mathrm{G}} (\lambda a) = \left(\exp_ {\mathrm{G}} \circ \mathrm{L} (i)\right) (\lambda a) = (i \circ \exp_ {\mathrm{H}}) (\lambda a) \in \mathrm{H}.
\]

因此 (i) \(\Rightarrow\) (ii)。H 有可数基时的逆向蕴含由命题 3 得出。

推论 3. 设 G 是李群，\(\rho\) 是 G 的解析线性表示，\(x \in L(G)\)，且 \(g \in G\)。

\begin{enumerate}
\def\labelenumi{(\roman{enumi})}
\item
  \(\rho(\exp x) = \exp L(\rho)x;\)
\item
  \(\mathrm{Ad}(\exp x) = \exp \mathrm{ad} x;\)
\item
  \(g(\exp x)g^{-1} = \exp(\operatorname{Ad} g.x).\)
\end{enumerate}

论证与 § 4, no. 4, Corollaries 2 and 3 to Proposition 8 中的相同。

推论 4. 设 G 是有限维连通李群。

\begin{enumerate}
\def\labelenumi{(\roman{enumi})}
\item
  \(\text{Int}(L(G)) = \text{Ad}(G).\)
\item
  令 \(\mathbf{Z}\) 为 \(\mathbf{G}\) 的中心。则 \(\mathbf{Z}\) 是 \(\mathbf{G}\) 的李子群，其李代数是 \(\mathbf{L}(\mathbf{G})\) 的中心。从 \(\mathbf{G} / \mathbf{Z}\) 到 \(\operatorname{Int} \mathbf{L}(\mathbf{G})\) 的映射由 \(g \mapsto \operatorname{Ad} g\) 通过商得到，且是李群同构。
\end{enumerate}

断言 (i) 由推论 3 (ii) 以及定义 2 后的注记得出。令 \(g \in G\)。要使 \(\operatorname{Ad} g = \operatorname{Id}_{\operatorname{L(G)}}\)，充要条件是 \(\operatorname{Int} g\) 与 \(\operatorname{Id}_G\) 在 \(e\) 的某个邻域上重合（§ 4, no. 1, Theorem 1 (ii)），因而在整个 \(G\) 上重合；换言之，充要条件是 \(g \in Z\)。于是 (ii) 由 Proposition 1 的 Corollary 1 得出。

定义 3. 设 G 是有限维连通李群。李群 \(\operatorname{Int}(\mathbf{L}(\mathbf{G})) = \operatorname{Ad}(\mathbf{G})\) 称为 G 的伴随群。

命题 11. 设 \(\mathbf{G}\) 是连通交换李群。

\begin{enumerate}
\def\labelenumi{(\roman{enumi})}
\item
  exp 是从加法李群 L(G) 到 G 的 étale 态射。
\item
  若 \(\mathbf{K} = \mathbf{R}\) 且 \(\mathbf{G}\) 有限维，则 \(\mathbf{G}\) 同构于形如 \(\mathbf{R}^p \times \mathbf{T}^q\) 的李群（\(p, q\) 为整数 \(\geqslant 0\)）。
\end{enumerate}

由 Hausdorff 公式，有 \((\exp x)(\exp y) = \exp (x + y)\)，当 \(x,y\) 充分接近 0 时成立；由解析延拓可知，该等式对所有 \(x,y\) 属于 \(\mathrm{L}(G)\) 时成立。因此 exp 是群同态，并且是 étale 的，因为

\[
\mathrm{T} _ {e} (\exp) = \mathrm{Id} _ {\mathrm{L(G)}}.
\]

从而得到 (i)。断言 (ii) 由 (i) 和 General Topology, Chapter VII, § 1, Proposition 9 得出。

命题 12. 设 G 是李群，L = L(G)。对于所有 \(x \in L\)，将 \(\mathrm{T}_{x}(\mathrm{L})\) 与 L 识别，从而 exp 在 x 处的右微分 \(\varpi(x)\) 是从 L 到 L 的线性映射。对于所有 \(x \in L\)，

\[
\varpi (x) = \sum_ {n \geqslant 0} \frac {1}{(n + 1) !} (\mathrm{ad} x) ^ {n}.
\]

两边都是 \(x\) 的解析函数，并且当 \(x\) 充分接近 0 时相等（§ 4, no. 3, Proposition 6）。

注记。\(\varpi(x) . (\mathrm{ad} x) = \exp \mathrm{ad} x - 1\)。我们滥用记号写成

\[
\varpi (x) = \frac {\exp \operatorname{ad} x - 1}{\operatorname{ad} x}.
\]

推论。设 \(\mathbf{G}\) 是复李群，\(x \in \mathrm{L}(\mathbf{G})\)。在 \(x\) 处，\(\exp_{\mathbf{G}}\) 的切映射的核为 \(\bigoplus_{n \in \mathbf{Z}\setminus\{0\}} \operatorname{Ker}(\operatorname{ad} x - 2i\pi n)\)。

整函数 \(z \mapsto \sum_{n \geqslant 0} \frac{1}{(n + 1)!} z^n\) 等于 \(\frac{e^z - 1}{z}\)（当 \(z \neq 0\) 时），并且具有

零点 \(2\pi i\mathbf{Z}\setminus\{0\}\) 中的各点，且这些都是单零点。于是推论由命题 12 和以下引理得出：

引理 2. 设 E 是复 Banach 空间，u 是 \(\mathcal{L}(\mathrm{E})\) 的元素，S 是 \(u\) 在 \(\mathcal{L}(\mathrm{E})\) 中的谱（Spectral Theories, Chapter I, §1, no. 2），\(f\) 是 S 的一个开邻域 \(\Omega\) 上的全纯复函数。假设 \(f\) 在 \(\Omega\) 中只有有限个互异零点 \(z_{1},\ldots ,z_{n}\)，其重数分别为 \(h_1,\dots,h_n\)。则 \(\operatorname {Ker}f(u)\) 是各 \(\operatorname {Ker}(u - z_i)^{h_i}\)（\(1\leqslant i\leqslant n\)）的直和。

（关于 \(f(u)\) 的定义，见 Spectral Theories, Chapter I, § 4, no. 8。）

存在处处非零的全纯函数 \(g\)，它定义在 \(\Omega\) 上，使得 \(f(z) = (z - z_1)^{h_1} \ldots (z - z_n)^{h_n} g(z)\)。于是 \(g(u)g^{-1}(u) = g^{-1}(u)g(u) = 1\)，从而 \(\operatorname{Ker}f(u) = \operatorname{Ker}\prod_{i=1}^{n}(u - z_i)^{h_i}\)。将 \(\operatorname{Ker}f(u)\) 看作 \(\mathbf{C}[\mathbf{X}]\)-模，其外作用律为 \((h,x)\mapsto h(u)x\)，其中 \(h\in \mathbf{C}[\mathbf{X}]\)、\(x\in \operatorname{Ker}f(u)\)，可见 \(\operatorname{Ker} f(u)\) 是各 \(\operatorname{Ker}(u - z_i)^{h_i}\) 的直和，应用 Algebra, Chapter VII, § 2, no. 1, Proposition 1 即得。

